%Diese LaTeX-Vorlage f�r Praktikumsprotokolle in Form einer Ver�ffentlichung f�r das 
%F-Praktikum wurde erstellt von Andreas Nuber EP II, Uni W�rzburg.
%Es wurde das Koma-Skript verwendet und sollte somit installiert sein
%desweiteren sollten alle packages installiert sein, die mit \usepackage{} 
%eingebunden werden.
%Zum Testen dieser Vorlage wurde MikTex verwendet sowie TeXnicCenter als Editor
%beim Compilieren waren 0 Fehler, 1 Warnung und 3 zu volle/leere Boxen. Das ist ok :)

%F�r das Erstellen einfach den sinnfreien Text, der zum ausf�llen genommen wurde 
%ersetzen. Was sonst noch ver�ndert werden sollte steht in den Kommentaren!


\documentclass[a4paper,10pt,twocolumn]{scrartcl} %Koma-Skript-�quivalent zu "article"

\usepackage{german}            %macht deutsche �berschriften
\usepackage{amsmath}           %macht
\usepackage{amsfonts}          %       Mathe
\usepackage{amssymb}           %              m�chtiger
\usepackage{graphicx}          %erlaubt Graphiken einzubinden (.eps f�r dvi und ps sowie .jpg f�r pdf)
\usepackage[T1]{fontenc}       %Zeichenbelegung der verwendeten Schrift
\usepackage{ae}                %macht sch�neres �
\usepackage{typearea}	         %erm�glicht �nderung des Seitenspiegels
\usepackage{scrlayer-scrpage}          %erm�glicht �nderung der Kopf-/Fu�zeile
\usepackage{lastpage}         %l�sst auf die Seienanzahl zugreifen
\usepackage[margin=10pt,font=small,labelfont=bf]{caption} %macht die Bildbeschriftungen richtig
\usepackage{xcolor}			  % macht Farbe
\usepackage{subfig}
\usepackage{hyperref}
\usepackage{float}
\usepackage{comment}

\renewcommand{\figurename}{Abb.}

\pagestyle{scrheadings}        %sagt Koma-Skript, dass selbstdefiniers Kopfzeilen verwendet werden
\typearea{16}                  %stellt Seitenspiegel ein
\columnsep25pt								 %definiert Breite zwischen den zwei Spalten von \twocolumns

\renewcommand{\pnumfont}{%     %�ndert die Schriftart der Seitennummerierung
\normalfont\rmfamily\slshape}  %�ndert die Schriftart der Seitennummerierung 



\begin{document}
\cfoot{\thepage /\pageref{LastPage}}     %macht die Seitennumerierung der Fom 2/5 (ausser auf der Titelseite)

\twocolumn[{\csname @twocolumnfalse\endcsname                %erlaubt "Abstrakt" �ber beide Spalten
\titlehead{                                                  %Kopfzeile
	\begin{tabular*}{\textwidth}[]{@{\extracolsep{\fill}}lr}   %Kopfzeile
	Betreuer: Prof. Dr. V. Hinkov & \today\\                          %Kopfzeile      hier den Betreuer eintragen!!!
	\end{tabular*}                                             %Kopfzeile
	}
\title{Plancksches Wirkungsquantum und Boltzmann-Konstante}  %Titel der Versuchs
\author{Moritz Langer \and Hannes Wikler}                     %Namen der Studenten
\date{}			                                               %ben�tigt um automatisches Datum auszuschalten
\maketitle                                                      %erzeugt Titelseite
\vspace{-8ex}                                                   %verringert Abstand zur �berschrift
\begin{abstract}                                                %Beginn des Abstracts
\vspace{1em}
\noindent Versuchsdurchführung: 20. November 2025\\% Datum ändern!
\noindent Protokollabgabe: 4. Dezember 2025% Datum ändern!
\vspace{1em}

	\noindent Das Plancksche Wirkungsquantum wurde durch den Photoeffekt als $\textit{h} = (3.32 \pm 0.12)\cdot 10^{-34} \, J s$ bestimmt - damit
	ist es nicht mit dem Literaturwert von $h_{lit} = 6.62607015 \cdot 10^{-34} \, J s$ \cite{PlanckKonstante} vereinbar. \\
	Zusätzlich wurde die Boltzmannkonstante durch Messung des Diffusionsstroms eines npn-Leistungstransistors bei verschiedenen Temperaturen auf
	$k_B = (1.4291 \pm 0.0087) \cdot 10^{-23} J/K$ bestimmt - damit ist auch dieser Wert ist nicht mit dem Literaturwert von 
	$k_{B,lit} = 1.380649 \cdot 10^{-23} J/K$ \cite{BoltzmannKonstante} vereinbar, liegt
\end{abstract}
}]

\section{Einleitung}
\section{Theorie}
\section{Experiment}
\section{Zusammnfassung}
\section{Anhang}
    \color{blue}{HIER NOCH RECHNUNG UND HERLEITUNG REIN}
    Die Isotherme für C02 dargestellt:
	
\begin{thebibliography}{}    %so wird das Literaturverzeicnis erstellt
\bibitem{BipBop}
	\url{https://www.physik.uni-wuerzburg.de/fileadmin/11000000/03_Studium/02_Bachelor/Grundpraktikum/_imported/fileadmin/11016800/Modulbeschreibungen/C-Modul/Versuch_41_ueberarbeitet.pdf}, besucht am 02.12.2025 um 20.13 Uhr
\end{thebibliography}

\end{document}